% !TEX program = xelatex
% !TEX root = ../../TFM.tex
\documentclass[../../TFM.tex]{subfiles}

\begin{document}

% =====================================
% SECCIÓN DE METODOLOGÍA
% =====================================


	\chapter{Metodología y estrategia empírica}
	\label{chap:metodologia}
	\thesisstartchap{metodologia}


	\section{Tipo de estudio y enfoque}
	\label{sec:tipo-de-estudio}

	El estudio adopta un enfoque \emph{cuantitativo} con alcance
	\emph{explicativo}, orientado a contrastar empíricamente las hipótesis
	formuladas sobre la relación entre rentas de recursos naturales,
	condiciones institucionales y transformación estructural en economías
	dependientes de recursos. El diseño de investigación es
	\emph{no experimental} y \emph{longitudinal} (panel), y se basa en el
	uso de datos secundarios comparables internacionalmente para un conjunto
	de países. La cobertura histórica de los insumos varía según la fuente:
	Penn World Table ofrece series desde 1950; el índice CTOT del Fondo
	Monetario Internacional, desde 1962; los \textit{World Development
	Indicators}, desde 1980; la información comercial del
	\textit{Atlas of Economic Complexity} y los años promedio de escolaridad
	del PNUD, desde 1990; y los
	\textit{Worldwide Governance Indicators}, desde 1996. El período común de
	análisis econométrico se fija en 1996--2021: comienza con la disponibilidad
	de los indicadores institucionales y termina en el último año para el cual
	se publican conjuntamente los cuatro componentes requeridos para construir
	$RENTS$. El intervalo 1990--1995 se utiliza exclusivamente como período base
	para clasificar la dependencia externa de recursos naturales mediante el
	indicador $DRES$.

	En este marco, la estrategia empírica combina análisis descriptivo y dos
	familias de estimaciones en datos de panel con efectos fijos bidireccionales.
	La familia principal utiliza el índice de complejidad económica ($ECI$) como
	variable dependiente, mientras que la familia complementaria utiliza
	$DIVX=1-HHI$. Ambas incorporan el mismo conjunto de canales teóricos y la
	interacción entre rentas extractivas y calidad institucional; la única
	diferencia entre ellas es la exclusión de $HHI$ cuando $DIVX$ constituye la
	variable dependiente. Dentro de cada familia, M1 y M2 son especificaciones
	temáticas paralelas y M3 es la especificación completa principal. Las dos
	ecuaciones M3 definen los resultados sustantivos del estudio, mientras que la
	desagregación por tipo de recurso y las pruebas de estabilidad se utilizan para
	evaluar la sensibilidad de la evidencia principal.

	Dado el carácter observacional del diseño y la posible endogeneidad
	entre algunas variables centrales, los resultados se interpretarán como
	asociaciones condicionales y no como estimaciones causales estrictas. Los
	efectos fijos por país controlan la heterogeneidad inobservable constante
	en el tiempo y los efectos temporales absorben shocks globales comunes,
	aunque estas decisiones no eliminan por completo las limitaciones propias
	de un diseño observacional.

	La estrategia metodológica adoptada se deriva directamente de los
	mecanismos productivos, institucionales y macroeconómicos discutidos en
	el \autoref{chap:marco-teorico}, y tiene como objetivo contrastar
	empíricamente las hipótesis formuladas en el
	\autoref{chap:hipotesis}.
	\section{Unidad de análisis, población y muestra}
	\label{sec:unidad-analisis}

	La unidad de análisis es el \emph{país--año}. La población de referencia
	comprende economías para las cuales existe información comparable sobre
	dependencia extractiva, transformación estructural, condiciones
	institucionales y variables macroeconómicas relevantes durante el período
	principal de estimación 1996--2021. A partir de este universo, la base
	empírica se organiza como un panel internacional, potencialmente
	desbalanceado,\footnote{Un panel desbalanceado permite que no todos los países cuenten con observaciones para todos los años, siempre que la ausencia de datos sea tratada de forma consistente con la estrategia econométrica y no responda sistemáticamente al resultado que se busca explicar.} en función de la
	disponibilidad y consistencia temporal de los datos.

	La muestra analítica principal se concentra en economías dependientes de
	recursos naturales, identificadas a partir de criterios explícitos de
	dependencia externa que se desarrollan en la sección siguiente. Esta
	decisión responde al objetivo central del trabajo, que no consiste en
	comparar cualquier tipo de economías, sino en examinar las diferencias en
	las trayectorias de transformación estructural dentro del conjunto de
	países cuya inserción internacional está fuertemente condicionada por los
	recursos naturales.

	De manera complementaria, en algunas etapas descriptivas podrán
	considerarse economías con baja dependencia de recursos como grupo de
	referencia, con el fin de contextualizar los patrones observados en las
	economías extractivas. No obstante, el núcleo del análisis econométrico
	recae sobre la muestra de países dependientes de recursos, seleccionada
	según criterios de disponibilidad, comparabilidad internacional y
	relevancia analítica. Esta selección fija una grilla maestra de 55 países y
	26 años, equivalente a 1.430 observaciones país--año. La cantidad de
	observaciones y países efectivamente utilizada puede variar entre
	especificaciones por la intersección de las variables requeridas, pero no
	se redefine la muestra de países ni se reemplazan sus integrantes.

	\section{Criterios de selección de la muestra}
	\label{sec:criterios-seleccion-muestra}

	La construcción de la muestra sigue un procedimiento secuencial orientado
	a identificar economías con elevada dependencia externa de recursos
	naturales no renovables del subsuelo. Aunque la unidad de análisis del
	estudio es el país--año, la pertenencia a cada muestra se define a nivel de
	país durante un período anterior a la estimación, con el propósito de evitar
	que la clasificación quede afectada por trayectorias posteriores de
	transformación estructural. La
	\autoref{fig:sample_selection_pipeline} resume este procedimiento.

	El universo inicial corresponde a las 233 economías del catálogo comercial
	del \textit{Atlas of Economic Complexity} para la clasificación SITC
	Revisión 2. Se exige que las exportaciones de mercancías sean positivas en cada uno
	de los seis años comprendidos entre 1990 y 1995, condición que cumplen 185
	economías. Posteriormente se restringe el universo a 159 países, excluyendo
	territorios dependientes, regiones administrativas especiales, entidades
	históricas disueltas y códigos comerciales sin un país identificable. Esta
	delimitación conserva países con independencia del tamaño de su población y
	del valor absoluto de sus exportaciones.

	Los umbrales de un millón de habitantes y de mil millones de dólares de
	exportaciones no se utilizan como exclusiones de la muestra. Ambas
	variables se conservan únicamente como diagnósticos descriptivos, pues su
	aplicación mecánica excluiría economías extractivas pequeñas o de bajos
	ingresos y podría introducir un sesgo de selección asociado al tamaño del
	país, en lugar de identificar su grado de dependencia externa de recursos.

	Sobre este conjunto de economías se identifica el grupo de países con
	elevada dependencia externa de recursos naturales no renovables del subsuelo. Para ello se construye
	el indicador $DRES_{it}$, definido como la participación de las
	exportaciones de estos recursos en el total de exportaciones de mercancías del
	país, tal como se presenta en la \autoref{eq:selection_dependence}:

	\begin{equation}
		DRES_{it} =
		\frac{X^{\text{res}}_{it}}{X_{it}}
		\label{eq:selection_dependence}
	\end{equation}

	donde $DRES_{it}$ representa el grado de dependencia externa de recursos
	naturales no renovables del subsuelo del país $i$ en el año $t$;
	$X^{\text{res}}_{it}$ denota el valor de las exportaciones de combustibles
	minerales, minerales crudos, minerales metalíferos, metales no ferrosos,
	piedras preciosas y oro no monetario del país $i$ en el año $t$; y
	$X_{it}$ corresponde al valor total de las exportaciones de mercancías del
	país $i$ en el año $t$.

	La construcción utiliza exportaciones a cuatro dígitos de la clasificación
	SITC Revisión 2. El numerador reúne la sección 3, las divisiones 27, 28 y 68,
	los códigos 6672 y 6673, y el código 9710. Estos últimos incorporan
	explícitamente diamantes y piedras preciosas naturales, y oro no monetario,
	recursos del subsuelo que no están incluidos en el indicador estándar de
	minerales y metales de los WDI. Dentro de los grupos amplios se excluyen el
	código 2711, correspondiente a fertilizantes de origen animal o vegetal; el
	código 3510, correspondiente a electricidad; los códigos 2786, 2820, 2881 y
	2882, que corresponden a escorias, cenizas, residuos y chatarra; las perlas
	del código 6671; y las piedras sintéticas o reconstruidas del código 6674. La
	definición excluye así bienes agrícolas, forestales, pesqueros y los demás
	recursos naturales renovables, además de electricidad, materiales recuperados
	y productos sintéticos.

	El denominador incorpora todas las mercancías con código SITC y conserva la
	categoría \texttt{XXXX}, utilizada por el Atlas para comercio de bienes cuyo
	producto no pudo clasificarse. Mantener esta categoría evita elevar
	mecánicamente $DRES$ al retirar exportaciones de mercancías que sí forman parte
	del total comercial. En cambio, se excluyen expresamente las categorías de
	servicios financieros, servicios empresariales, transporte, viajes y servicios
	no especificados. Esta distinción resulta necesaria porque los archivos del
	Atlas reúnen comercio de bienes y servicios, mientras que $DRES$ se define
	exclusivamente sobre exportaciones de mercancías.

	La clasificación de economías dependientes se realiza utilizando el
	promedio de $DRES_{it}$ durante el período base 1990--1995. Esta decisión
	permite definir la pertenencia a la muestra con anterioridad a buena
	parte de las trayectorias posteriores de diversificación o cambio
	estructural, reduciendo así problemas de simultaneidad entre la
	condición extractiva y los resultados que se pretende explicar. Sobre esa
	base, se define un umbral $\theta$ que representa la participación mínima
	de exportaciones de recursos naturales en el total exportado requerida
	para clasificar una economía como dependiente de recursos. La estimación
	de referencia utiliza $\theta=20\%$, umbral empleado por el Fondo Monetario
	Internacional para identificar economías en las que los recursos agotables
	representan una proporción significativa de las exportaciones
	\parencite{IMF2012}. Este valor se adopta como referencia externa y no como
	una réplica exacta de la clasificación del FMI, originalmente construida para
	otro período y un universo específico de economías en desarrollo. Con los
	datos del período base, esta regla produce una muestra fija de 55 países. La
	composición nominal de la muestra y el valor promedio de $DRES$ de cada país
	se presentan en la \autoref{tab:muestra_dres20}.

	Como validación de la construcción de $DRES$, se examinan decisiones sobre
	el ámbito de productos y la agregación temporal. En primer lugar, se
	construye una variante que excluye oro no monetario y piedras preciosas, con el
	fin de reducir la posible influencia de reexportaciones, y una definición
	restrictiva que conserva solamente productos extractivos ubicados aguas arriba.
	En segundo lugar, el promedio simple de los seis valores anuales se contrasta
	con la mediana, el cociente entre exportaciones extractivas acumuladas y
	exportaciones acumuladas de mercancías, y ejercicios que retiran un año del
	período base cada vez. El cociente de exportaciones acumuladas conserva los 55
	países de la muestra principal del 20\%, mientras que la mediana conserva 54;
	además, 53
	permanecen clasificados al retirar cualquiera de los seis años, mientras que
	las variantes sin oro y piedras preciosas y estrictamente aguas arriba contienen
	49 y 48 países, respectivamente. Estos cálculos se utilizan exclusivamente
	para verificar la clasificación y no definen muestras econométricas
	adicionales: todas las estimaciones emplean el umbral fijo de 20\%.

	Como validación concurrente, la construcción se compara con la suma de los
	indicadores WDI de exportaciones de combustibles y de minerales y metales. En
	las 495 observaciones país--año comparables de 123 países, la variante del Atlas
	que reproduce esas categorías alcanza una correlación de 0,974 y un error
	absoluto mediano de 0,58 puntos porcentuales frente a WDI. Debido a la elevada
	presencia de valores faltantes en estos indicadores durante 1990--1995, WDI se
	utiliza como fuente de validación y no como insumo principal de la selección.

	Este criterio de selección responde a que el interés del estudio se
	centra en la dependencia externa asociada a la estructura exportadora, y
	no en la mera disponibilidad geológica o física de recursos. En
	consecuencia, economías que poseen recursos naturales pero los destinan
	principalmente al consumo interno no se clasifican como dependientes de
	recursos en este trabajo. Al basarse en exportaciones brutas, la medida tampoco
	distingue perfectamente entre extracción doméstica, procesamiento y
	reexportación; por ello, $DRES$ debe interpretarse como un indicador de
	dependencia y especialización exportadora, no como una medida directa de
	dotación geológica o producción física. Asimismo, conviene subrayar que $DRES$ se
	emplea exclusivamente como criterio de selección muestral, mientras que
	la variable principal de interés en las estimaciones econométricas es
	$RENTS_{it}$, que captura la intensidad macroeconómica de las rentas
	extractivas como porcentaje del PIB.

	La elección de una medida basada en exportaciones sigue a la literatura
	que vincula la dependencia de recursos con la especialización externa y
	con los límites a la diversificación productiva. No obstante, se reconoce
	que esta aproximación puede introducir sesgos hacia economías de menor
	ingreso, en la medida en que países con menor desarrollo productivo
	tienden a concentrar sus exportaciones en recursos naturales
	\parencite{Brunnschweiler2008}. Frente a esta limitación, la literatura ha
	propuesto medidas alternativas basadas en rentas de recursos, ya sea como
	participación en el PIB o en términos per cápita. Sin embargo, dichas
	medidas capturan una dimensión distinta del fenómeno: mientras $DRES$
	refleja dependencia externa y especialización exportadora, $RENTS$
	aproxima la importancia macroeconómica de la actividad extractiva. Por
	ello, ambas medidas se utilizan en este capítulo con funciones analíticas
	diferenciadas.

	La disponibilidad de PIB y de las demás variables del modelo no interviene
	en el cálculo de DRES ni en la clasificación inicial. La muestra efectiva de
	cada especificación corresponde posteriormente a la intersección entre la
	lista definida por el umbral y la información disponible para las variables
	econométricas. Por ello, el número de países y observaciones puede disminuir
	en las estimaciones sin modificar la definición previa de dependencia
	extractiva.

	\begin{figure}[htbp]
		\centering

		\begin{tikzpicture}[
			node distance=1.2cm,
			every node/.style={font=\small},
			arrow/.style={->, thick}
			]

			% Cajas con anchos decrecientes
			\node[draw, rectangle, minimum width=7.2cm, minimum height=0.8cm, align=center] (all)
			{Catálogo Atlas SITC: 233 economías};

			\node[below=of all] (text1)
			{Mercancías positivas en cada año de 1990--1995};

			\node[draw, rectangle, minimum width=6.4cm, minimum height=0.8cm, align=center, below=of text1] (step1)
			{185 economías};

			\node[below=of step1] (text2)
			{Ámbito nacional del estudio};

			\node[draw, rectangle, minimum width=5.6cm, minimum height=0.8cm, align=center, below=of text2] (step2)
			{159 países};

			\node[below=of step2, align=center] (text3)
			{$DRES$ promedio 1990--1995 $\geq \theta$};

			\node[draw, rectangle, minimum width=4.8cm, minimum height=0.8cm, align=center, below=of text3] (step3)
			{$\theta=20\%$: 55 países};

			\node[below=of step3, align=center] (text4)
			{Muestra fija de dependencia extractiva};

			\node[draw, rectangle, minimum width=4.0cm, minimum height=0.8cm, align=center, below=of text4] (step4)
			{55 países seleccionados};

			\node[below=of step4, align=center] (text5)
			{Datos completos \\ para todas las variables};

			\node[draw, rectangle, minimum width=3.2cm, minimum height=0.8cm, align=center, below=of text5] (final)
			{Muestra efectiva \\ del panel};

			% Flechas
			\draw[arrow] (all) -- (text1);
			\draw[arrow] (text1) -- (step1);
			\draw[arrow] (step1) -- (text2);
			\draw[arrow] (text2) -- (step2);
			\draw[arrow] (step2) -- (text3);
			\draw[arrow] (text3) -- (step3);
			\draw[arrow] (step3) -- (text4);
			\draw[arrow] (text4) -- (step4);
			\draw[arrow] (step4) -- (text5);
			\draw[arrow] (text5) -- (final);

			% Embudo
			\draw
			([xshift=-1cm]all.north west) --
			([xshift=1cm]all.north east);

			\draw
			([xshift=-1cm]all.north west) --
			([xshift=-0.6cm]final.south west);

			\draw
			([xshift=1cm]all.north east) --
			([xshift=0.6cm]final.south east);

		\end{tikzpicture}

		\caption[Selección de la muestra]{Procedimiento reproducible de selección de países utilizado para construir la muestra analítica del estudio. La regla única aplica $\theta=20\%$ sobre la dependencia externa de recursos, medida mediante $DRES$. Los conteos corresponden a los resultados efectivos del procesamiento del período 1990--1995.}

		\label{fig:sample_selection_pipeline}
		\caption*{\footnotesize \textit{Fuente: Elaboración propia con base en el Atlas of Economic Complexity del Harvard Growth Lab.}}

	\end{figure}

	\section{Heterogeneidad estructural del extractivismo}
	\label{sec:heterogeneidad-estructural}

	Aunque la muestra analítica principal se restringe a economías
	dependientes de recursos naturales, este conjunto de países no constituye
	un bloque homogéneo. Por el contrario, dentro de las economías
	extractivas existen diferencias relevantes asociadas al tipo de recurso
	predominante, al grado de especialización exportadora y al nivel de
	desarrollo económico. Estas diferencias reflejan trayectorias históricas,
	institucionales y productivas diversas, y resultan importantes para
	interpretar los patrones observados en la transformación estructural.

	No obstante, el objetivo central del estudio no es construir una tipología
	exhaustiva del extractivismo ni comparar formalmente todos sus subtipos,
	sino identificar regularidades generales dentro del conjunto de economías
	seleccionadas por su dependencia externa de recursos. Por esta razón, la
	estimación econométrica principal se realiza sobre la muestra completa de
	países dependientes de recursos, preservando comparabilidad
	internacional, consistencia conceptual y tamaño muestral suficiente.

	En este marco, la heterogeneidad estructural se incorpora principalmente
	con fines descriptivos y de interpretación. En primer lugar, los países
	pueden clasificarse según el tipo de recurso predominante en su estructura
	exportadora, distinguiendo entre economías dependientes de petróleo, gas y
	minería. Esta distinción permite
	identificar diferencias básicas en la composición de la canasta
	exportadora y en los mecanismos productivos asociados a cada perfil
	extractivo.

	En segundo lugar, se considera el grado de dependencia externa de
	recursos, medido a partir de la participación de las exportaciones de
	recursos naturales en el total exportado. Con fines descriptivos, esta
	dimensión puede ordenarse en categorías de dependencia moderada, alta y
	muy alta, lo que permite observar si los patrones de transformación
	estructural difieren según la intensidad de la especialización
	exportadora. Esta clasificación no sustituye el criterio principal de
	selección muestral basado en el umbral $\theta$, sino que lo complementa
	mediante una caracterización más fina de la muestra seleccionada.

	Finalmente, se incorpora una clasificación por nivel de ingreso de acuerdo
	con la tipología del Banco Mundial ---\textit{ingreso bajo (LICs)},
	\textit{ingreso medio-bajo (LMICs)}, \textit{ingreso medio-alto (UMICs)}
	y \textit{ingreso alto (HICs)}---, siguiendo el enfoque de
	\textcite{Anne2021}. Con el fin de evitar cambios en la composición de los
	grupos a lo largo del tiempo, se adopta una clasificación fija basada en
	el nivel de ingreso correspondiente al año 2010. Esta dimensión permite
	contextualizar la heterogeneidad de las economías extractivas en función
	de sus capacidades iniciales, su estructura productiva y su margen de
	acción institucional.

	En conjunto, estas dimensiones no redefinen la muestra ni alteran el
	criterio principal de inclusión, pero sí ofrecen un marco útil para
	describir la variedad interna de las economías dependientes de recursos y
	para interpretar, en etapas posteriores, posibles diferencias en sus
	trayectorias de transformación estructural.
	\section{Variables del estudio y definición operacional}
	\label{sec:variables-estudio}

La definición operativa de recursos naturales adoptada en este trabajo se
alinea con la evolución reciente de la literatura sobre la denominada
``maldición de los recursos''. Mientras estudios iniciales como
\textcite{Sachs1995, Sachs1997} incluían dentro de los recursos naturales
a un conjunto amplio de productos primarios ---incluyendo bienes
agrícolas---, trabajos posteriores han restringido esta noción a
recursos de tipo extractivo, particularmente hidrocarburos y minerales
\parencite{Isham2005}. Esta delimitación resulta más adecuada para el
problema de investigación, ya que los recursos extractivos dependen en
mayor medida de la dotación geológica del país y presentan rasgos
económicos específicos ---como agotabilidad, concentración geográfica y
facilidad relativa de apropiación--- que los vinculan estrechamente con
dinámicas rentistas de carácter institucional, fiscal y macroeconómico
\parencite{Boschini2007, Lashitew2021}.

En consonancia con este criterio, el estudio se centra en economías
históricamente dependientes de recursos naturales no renovables,
particularmente economías dependientes de recursos del subsuelo como petróleo,
gas, carbón, minerales y metales. Esta definición excluye expresamente los bienes
agrícolas, forestales y pesqueros, así como los demás recursos naturales renovables.
No obstante, la
dependencia de recursos y la importancia macroeconómica de la actividad
extractiva no se tratan como conceptos equivalentes. La primera se
aproxima mediante $DRES_{it}$ y se utiliza exclusivamente como criterio
de selección muestral, mientras que la segunda se representa mediante
$RENTS_{it}$, que captura el peso de las rentas extractivas del subsuelo
como porcentaje del PIB y constituye la variable principal de interés en
las estimaciones econométricas. De este modo, se distingue entre la
condición inicial de dependencia externa y la intensidad económica de las
rentas extractivas dentro de la muestra seleccionada.

La variable dependiente central del estudio es la transformación
estructural. La literatura ha tendido a aproximarse a este fenómeno en
economías dependientes de recursos mediante la expansión del sector no
extractivo o mediante descomposiciones sectoriales entre actividades
vinculadas y no vinculadas a los recursos naturales
\parencite{Lashitew2021}. Si bien este enfoque resulta intuitivo, suele
ser sensible a problemas de clasificación, disponibilidad y
comparabilidad internacional de los datos. Por esta razón, este trabajo
adopta como aproximación principal el \textit{Economic Complexity Index}
($ECI_{it}$), entendido como una medida de las capacidades productivas
acumuladas y de la sofisticación de la estructura exportadora. En este
sentido, el ECI permite captar de forma más estructural la evolución de
las capacidades productivas, relativamente menos expuesta a fluctuaciones
coyunturales asociadas a los ciclos de precios de commodities.

Operativamente, se utiliza el campo \texttt{countryYear.eci} de la API
GraphQL oficial del \textit{Atlas of Economic Complexity}
\parencite{HarvardGrowthLab2026}. La versión vigente del Atlas representa la
intensidad de la especialización mediante
$M_{cp}=RCA_{cp}/(1+RCA_{cp})$, en lugar de reducirla a una presencia binaria.
Para mantener una definición homogénea, todas las observaciones se extraen de
una única captura reproducible de la API para 1995--2022 y no se combinan con
la antigua tabla de rankings ni con indicadores de otros proveedores. En el
período 1996--2021, esta fuente cubre los 55 países y las 1.430 observaciones
país--año de la grilla maestra, incluso cuando alguna economía no sea elegible
para aparecer en los rankings públicos del Atlas.

No obstante, dado que el ECI combina información sobre diversificación y
sofisticación exportadora, el análisis incorpora como resultado
complementario una medida más directa de diversificación de la canasta
exportadora. Esta variable, denotada como $DIVX_{it}=1-HHI_{it}$, permite
evaluar si los mecanismos asociados a mayores niveles de complejidad
económica también se relacionan con una menor concentración exportadora.
La medida se interpreta como complemento del ECI, no como sustituto de la
sofisticación productiva capturada por dicho índice.

En el plano explicativo, $RENTS_{it}$ se incorpora como uno de los componentes
del modelo, junto con variables que capturan distintas dimensiones de la
inserción extractiva y de sus canales de transmisión. Entre ellas se incluyen
las rentas reales por habitante provenientes del petróleo, el gas y el carbón
($OILPC_{it}$, $GASPC_{it}$ y $COALPC_{it}$), indicadores de especialización exportadora como la
participación de exportaciones primarias no energéticas ($PEXP_{it}$) y de combustibles
($FEXP_{it}$), y una medida de concentración exportadora ($HHI_{it}$).
Asimismo, la calidad institucional ($INST_{it}$) se incorpora como
variable moderadora para evaluar si la asociación entre las rentas extractivas
y la transformación estructural varía con el entorno institucional.
Junto con ello, se añaden variables macroeconómicas, fiscales, financieras
y de capacidades productivas, tales como el tipo de cambio real
($RER_{it}$), la volatilidad ($VOL_{it}$), el consumo final del gobierno
como porcentaje del PIB ($GOVCONS_{it}$), la profundidad financiera ($FIN_{it}$), el capital humano
($HUMCAP_{it}$), la innovación ($INNOV_{it}$), la conectividad digital
($NET_{it}$) y el logaritmo del PIB per cápita ($\log(GDPPC_{it})$),
utilizado como proxy del nivel de desarrollo económico.

La \autoref{tab:variables} presenta la descripción, definición
operacional y fuentes de las variables utilizadas en el análisis
empírico. La selección definitiva de indicadores responde a criterios de
	consistencia conceptual, comparabilidad internacional, cobertura temporal
	y disponibilidad de información para el período de estimación 1996--2021.
	\renewcommand{\arraystretch}{1.2}
	\setlength{\tabcolsep}{4pt}

	\begin{longtable}{|
			>{\raggedright\arraybackslash}p{0.22\textwidth}|
			>{\centering\arraybackslash\ttfamily}p{0.10\textwidth}|
			>{\raggedright\arraybackslash}p{0.14\textwidth}|
			>{\raggedright\arraybackslash}p{0.34\textwidth}|
			>{\raggedright\arraybackslash}p{0.15\textwidth}|}

		\caption[Variables y definiciones]{Descripción y definición operacional de las variables del modelo econométrico.}

		\label{tab:variables} \\
		\hline

		Variable
		& \multicolumn{1}{|>{\centering\arraybackslash}p{0.10\textwidth}|}{Símbolo}
		& Tipo
		& Definición operacional
		& Fuente \\

		\hline
		\endfirsthead

		\multicolumn{5}{c}{{\tablename\ \thetable{} -- continuación}} \\
		\hline

		Variable
		& \multicolumn{1}{|>{\centering\arraybackslash}p{0.10\textwidth}|}{Símbolo}
		& Tipo
		& Definición operacional
		& Fuente \\

		\hline
		\endhead

		\hline
		\multicolumn{5}{r}{\textit{Continúa en la página siguiente}} \\
		\endfoot

		\hline

		\multicolumn{5}{c}{%
			\footnotesize\textit{Fuente: Elaboración propia con base en fuentes internacionales comparables.}
		} \\[0.3em]

		\multicolumn{5}{p{\textwidth}}{%
			\textit{Nota: La transformación estructural se aproxima desde su dimensión externa mediante el Índice de Complejidad Económica (ECI) como variable dependiente principal y la diversificación exportadora ($DIVX=1-HHI$) como resultado complementario. Cada variable utiliza exclusivamente la definición y la transformación indicadas en esta tabla y en las ecuaciones siguientes.}
		} \\

		\endlastfoot

		% -----------------------------------------------------
		% VARIABLES DEPENDIENTES
		% -----------------------------------------------------

		Complejidad económica
		& ECI
		& Dependiente (externa)
		& Índice de Complejidad Económica que combina la diversidad de la canasta exportadora y la ubicuidad de sus productos. Se utiliza la serie homogénea \texttt{countryYear.eci} de la API oficial; valores mayores indican mayores capacidades productivas y sofisticación exportadora.
		& Atlas of Economic Complexity / Harvard Growth Lab (API GraphQL) \\
		\hline

		Diversificación exportadora
		& DIVX
		& Dependiente complementaria
		& Complemento del índice de concentración de Herfindahl--Hirschman, definido como $1-HHI$; valores mayores indican una canasta exportadora menos concentrada.
		& Atlas of Economic Complexity / elaboración propia \\
		\hline

		% -----------------------------------------------------
		% CANAL INSTITUCIONAL
		% -----------------------------------------------------

		Rentas extractivas del subsuelo (\% del PIB)
		& RENTS
		& Explicativa principal
		& Suma de las rentas de petróleo, gas natural, carbón y minerales como porcentaje del PIB. Excluye rentas forestales y otros recursos renovables.
		& World Development Indicators (Banco Mundial) \\
		\hline

		Calidad institucional
		& INST
		& Explicativa moderadora
		& Promedio simple de los estimadores WGI de control de la corrupción, estado de derecho y efectividad gubernamental, calculado únicamente cuando los tres están observados; proxy de capacidad institucional.
		& Worldwide Governance Indicators \\
		\hline

		Interacción entre rentas extractivas e instituciones
		& RENTS $\times$ INST
		& Interacción
		& Término de interacción que evalúa si la asociación entre la intensidad de las rentas extractivas y la transformación estructural varía con la calidad institucional.
		& Elaboración propia \\
		\hline

		% -----------------------------------------------------
		% RENTAS EXTRACTIVAS ESPECÍFICAS
		% -----------------------------------------------------

		Renta petrolera per cápita
		& OILPC
		& Explicativa de intensidad extractiva
		& Valor real de la renta petrolera por habitante, obtenido al aplicar la renta petrolera como porcentaje del PIB al PIB real y dividir el resultado por la población. En las estimaciones se utiliza $\ln(1+OILPC)$.
		& World Development Indicators / elaboración propia \\
		\hline

		Renta gasífera per cápita
		& GASPC
		& Explicativa de intensidad extractiva
		& Valor real de la renta de gas natural por habitante, obtenido al aplicar la renta gasífera como porcentaje del PIB al PIB real y dividir el resultado por la población. En las estimaciones se utiliza $\ln(1+GASPC)$.
		& World Development Indicators / elaboración propia \\
		\hline

		Renta carbonífera per cápita
		& COALPC
		& Explicativa de intensidad extractiva
		& Valor real de la renta del carbón por habitante, obtenido al aplicar la renta carbonífera como porcentaje del PIB al PIB real y dividir el resultado por la población. En las estimaciones se utiliza $\ln(1+COALPC)$.
		& World Development Indicators / elaboración propia \\
		\hline

		% -----------------------------------------------------
		% CANAL ESTRUCTURAL
		% -----------------------------------------------------

		Concentración de la estructura exportadora
		& HHI
		& Explicativa estructural condicionada
		& Índice de Herfindahl--Hirschman calculado como la suma de los cuadrados de las participaciones de los productos SITC Rev. 2 de cuatro dígitos en las exportaciones de mercancías clasificadas. Se excluyen los servicios y el residuo \texttt{XXXX}. Se omite como regresor cuando $DIVX=1-HHI$ es la variable dependiente.
		& Atlas of Economic Complexity / elaboración propia \\
		\hline

		Exportaciones primarias no energéticas (\% de exportaciones de mercancías)
		& PEXP
		& Explicativa estructural
		& Participación de las secciones 0, 1, 2 y 4 y de la división 68 de la SITC Rev. 2 en las exportaciones de mercancías; proxy de especialización externa en productos primarios no energéticos. El denominador incluye el residuo no clasificable XXXX y excluye los servicios.
		& Atlas of Economic Complexity / elaboración propia \\
		\hline

		Exportaciones de combustibles (\% de exportaciones de mercancías)
		& FEXP
		& Explicativa estructural
		& Participación de la sección 3 de la SITC Rev. 2 en las exportaciones de mercancías; indicador de especialización exportadora en recursos energéticos. Comparte con PEXP el denominador, que incluye XXXX y excluye los servicios, pero sus numeradores no se superponen.
		& Atlas of Economic Complexity / elaboración propia \\
		\hline

		% -----------------------------------------------------
		% CANAL MACROECONÓMICO
		% -----------------------------------------------------

		Volatilidad de los términos de intercambio de commodities
		& VOL
		& Explicativa macroeconómica
		& Desviación estándar móvil de cinco años de las variaciones logarítmicas anuales del índice de precios de exportaciones netas de commodities del FMI (\texttt{CEMPI\_CTOTNX\_TT}), utilizando exclusivamente ponderaciones fijas (\texttt{H\_FW\_IX}).
		& Fondo Monetario Internacional (CTOT) \\
		\hline

		Tipo de cambio real
		& RER
		& Explicativa macroeconómica
		& Logaritmo del nivel de precios del producto relativo a Estados Unidos (\texttt{pl\_gdpo}); un aumento representa una apreciación real.
		& Penn World Table 11.0 \\
		\hline

		% -----------------------------------------------------
		% CAPACIDADES PRODUCTIVAS
		% -----------------------------------------------------

		Capital humano
		& HUMCAP
		& Explicativa de capacidades
		& Exponencial de una función por tramos de los años promedio de escolaridad de adultos de 25 años o más; aplica los retornos educativos utilizados por PWT y aproxima capacidades productivas acumuladas. No se interpreta como años de escolaridad.
		& PNUD, \textit{Human Development Report 2025}; construcción propia \\
		\hline

		Innovación
		& INNOV
		& Explicativa de capacidades
		& Logaritmo de uno más los artículos científicos y técnicos por millón de habitantes; proxy de capacidades científico-tecnológicas y producción de conocimiento.
		& World Development Indicators (Banco Mundial) \\
		\hline

		Conectividad digital
		& NET
		& Explicativa de capacidades
		& Personas que utilizan internet como porcentaje de la población (\texttt{IT.NET.USER.ZS}); proxy de difusión tecnológica y acceso a información productiva.
		& World Development Indicators (Banco Mundial) \\
		\hline

		% -----------------------------------------------------
		% CONTROL
		% -----------------------------------------------------

		Logaritmo del PIB per cápita
		& \begin{tabular}[c]{@{}c@{}}LOG\_\\GDPPC\end{tabular}
		& Control
		& Logaritmo del Producto Interno Bruto per cápita en PPA y dólares internacionales constantes (\texttt{NY.GDP.PCAP.PP.KD}); se utiliza como proxy del nivel de desarrollo económico.
		& World Development Indicators (Banco Mundial) \\
		\hline

		% -----------------------------------------------------
		% CONTROLES ECONÓMICOS Y FINANCIEROS
		% -----------------------------------------------------

		Consumo final del gobierno (\% del PIB)
		& GOVCONS
		& Control fiscal
		& Gasto de consumo final del gobierno general como porcentaje del PIB (serie 16, componente 3); controla por el tamaño relativo del consumo público y no se interpreta como una medida directa de capacidad ni de prociclicidad fiscal.
		& \textit{National Accounts Main Aggregates} (Naciones Unidas) \\
		\hline

		Profundidad financiera
		& FIN
		& Explicativa financiera
		& Crédito doméstico al sector privado otorgado por bancos como porcentaje del PIB (\texttt{FD.AST.PRVT.GD.ZS}); proxy de la profundidad del crédito bancario privado, no de su acceso, eficiencia o calidad.
		& World Development Indicators (Banco Mundial) \\
		\hline

	\end{longtable}

	Las rentas extractivas se construyen únicamente cuando están observados sus
	cuatro componentes:
	\[
		RENTS_{it}
		=
		OILRENT_{it}
		+
		GASRENT_{it}
		+
		COALRENT_{it}
		+
		MINERALRENT_{it}.
	\]
	El panel conserva además dos desagregaciones para distinguir mecanismos sin
	incorporar petróleo y gas como regresores separados:
	\[
		\begin{aligned}
			RENTS^{OIL+GAS}_{it} &= OILRENT_{it}+GASRENT_{it}, \\
			RENTS^{MINING}_{it} &= COALRENT_{it}+MINERALRENT_{it}.
		\end{aligned}
	\]
	La especificación principal mantiene $RENTS_{it}$ como medida agregada; las
	dos desagregaciones quedan disponibles para examinar separadamente
	hidrocarburos y minería sin alterar el conjunto central de variables.
	Los valores faltantes no se sustituyen por cero. Los cuatro indicadores
	proceden de los \textit{World Development Indicators} y están expresados
	como porcentaje del PIB: rentas petroleras
	(\texttt{NY.GDP.PETR.RT.ZS}), gasíferas
	(\texttt{NY.GDP.NGAS.RT.ZS}), carboníferas
	(\texttt{NY.GDP.COAL.RT.ZS}) y minerales
	(\texttt{NY.GDP.MINR.RT.ZS}).

	El índice institucional se calcula como el promedio simple de los tres
	estimadores WGI, que se publican en una escala comparable:
	\[
		INST_{it}
		=
		\frac{
			CC_{it}
			+
			RL_{it}
			+
			GE_{it}
		}{3},
	\]
	donde $CC$ representa control de la corrupción, $RL$ estado de derecho y
	$GE$ efectividad gubernamental. El índice se construye únicamente cuando
	los tres componentes están disponibles; no se interpolan los vacíos
	estructurales de 1997, 1999 y 2001.

	Las medidas de rentas por habitante se construyen de forma homogénea:
	\[
		KPC_{it}
		=
		\frac{
			\left(R^{K}_{it}/100\right)Y^{real}_{it}
		}{
			POP_{it}
		},
		\qquad
		K\in\{OIL,GAS,COAL\},
	\]
	donde $R^{K}_{it}$ es la renta del recurso correspondiente como porcentaje
	del PIB, $Y^{real}_{it}$ es el PIB en dólares constantes y $POP_{it}$ es la
	población total. No se imputan componentes faltantes.

	La medida de volatilidad se construye a partir del
	\textit{Commodity Net Export Price Index} del FMI
	(\texttt{CEMPI\_CTOTNX\_TT}). Para cada país, primero se calcula la
	variación porcentual logarítmica anual del índice:
	\[
		g_{it}
		=
		100\left[
			\log(CTOT_{it})
			-
			\log(CTOT_{i,t-1})
		\right].
	\]
	A continuación, la volatilidad se define como la desviación estándar
	móvil de las cinco variaciones anuales más recientes:
	\[
		VOL_{it}
		=
		\operatorname{sd}\left(
			g_{i,t-4},
			\ldots,
			g_{it}
		\right).
	\]
	La variable utiliza exclusivamente el índice con ponderaciones fijas
	(\texttt{H\_FW\_IX}), de modo que los cambios de $VOL_{it}$ reflejen
	variaciones de los precios internacionales y no modificaciones
	contemporáneas de la composición comercial, especialmente relevante porque
	las variables dependientes también se construyen a partir de la estructura
	exportadora.

	La medida de capital humano se construye con los años promedio de escolaridad
	$s_{it}$ de la población de 25 años o más publicados por el PNUD
	\parencite{UNDP2025}. Para mantener la interpretación del índice de capital
	humano de PWT, se aplica su función de retornos educativos por tramos
	\parencite{GGDCHumanCapital2016}:
	\[
		\phi(s_{it})=
		\begin{cases}
			0{,}134s_{it}, & s_{it}\leq4,\\
			0{,}134(4)+0{,}101(s_{it}-4), & 4<s_{it}\leq8,\\
			0{,}134(4)+0{,}101(4)+0{,}068(s_{it}-8), & s_{it}>8,
		\end{cases}
		\qquad
		HUMCAP_{it}=\exp\!\left[\phi(s_{it})\right].
	\]
	Por tanto, $HUMCAP_{it}$ representa un índice de capacidades acumuladas y no
	los años de escolaridad observados directamente.

	La medida de innovación se construye como
	\[
		INNOV_{it}
		=
		\log\left(
			1
			+
			\frac{\text{artículos científicos y técnicos}_{it}}
			{\text{población}_{it}}
			\times 1.000.000
		\right).
	\]
	El numerador corresponde al indicador de artículos científicos y técnicos
	\texttt{IP.JRN.ARTC.SC} de los \textit{World Development Indicators}. Esta
	medida aproxima las capacidades científico-tecnológicas y la producción de
	conocimiento, más que la innovación comercial en sentido estricto.

	El tipo de cambio real y el control de desarrollo se construyen,
	respectivamente, como
	\[
		RER_{it}=\log(pl\_gdpo_{it})
		\qquad\text{y}\qquad
		LOG\_GDPPC_{it}=\log(GDPPC_{it}),
	\]
	donde \texttt{pl\_gdpo} proviene de Penn World Table 11.0 y $GDPPC$ es el
	PIB per cápita en PPA y dólares internacionales constantes
	(\texttt{NY.GDP.PCAP.PP.KD}) de los \textit{World Development Indicators}.


	\section{Fuentes y técnicas de recolección de datos}
	\label{sec:recoleccion-datos}

	El estudio utiliza exclusivamente datos secundarios recopilados y
	sistematizados por fuentes comparables a escala internacional: los
	\textit{World Development Indicators} y \textit{Worldwide Governance
	Indicators} del Banco Mundial, el \textit{Atlas of Economic Complexity} del
	Harvard Growth Lab, la base \textit{National Accounts Main Aggregates} de
	Naciones Unidas \parencite{UnitedNationsStatisticsDivision2026}, el Fondo
	Monetario Internacional, Penn World Table y la base de indicadores del
	\textit{Human Development Report 2025} del PNUD \parencite{UNDP2025},
	según la variable considerada y tal como se resume en la
	\autoref{tab:variables}. La selección responde a criterios de comparabilidad,
	cobertura temporal, transparencia metodológica y reproducibilidad.

	La preparación de los datos sigue una arquitectura de tres niveles. La capa
	\texttt{data/raw} conserva las capturas y extracciones de las fuentes sin
	aplicar las transformaciones econométricas. La capa
	\texttt{data/processed} implementa, mediante un script específico, la única
	definición activa de cada variable. Finalmente, el panel maestro integra
	estas salidas mediante uniones uno a uno por código ISO3 y año sobre la
	grilla fija de 55 países y 1996--2021. Esta separación permite rastrear cada
	valor desde el panel analítico hasta su insumo original y evita mezclar
	procesos de descarga, construcción e integración.

	La validación comprueba la unicidad de las 1.430 llaves país--año, la
	ausencia de duplicados, la equivalencia entre las salidas CSV y Stata, los
	rangos admisibles y las identidades de las variables derivadas, entre ellas
	$DIVX=1-HHI$ y $RENTS\times INST$. La grilla maestra es balanceada en sus
	identificadores, pero conserva explícitamente la disponibilidad desigual de
	las variables. La tesis no interpola observaciones ni completa vacíos
	mezclando fuentes. En $GOVCONS$, la serie activa procede íntegramente de
	Naciones Unidas y WDI se conserva como referencia de contraste. Los
	metadatos país--año clasifican 1.106 valores como directos, 297 como
	derivados por la fuente y 27 como mixtos por superposición de métodos; estos
	procedimientos no son imputaciones realizadas por la tesis. Los 16 ceros de
	Venezuela publicados por WDI entre 1996 y 2011 permanecen documentados en
	raw, pero no forman parte de la variable activa. En $FIN$, la definición
	activa corresponde íntegramente al crédito privado otorgado por bancos
	(\texttt{FD.AST.PRVT.GD.ZS}); la definición amplia
	(\texttt{FS.AST.PRVT.GD.ZS}) se conserva únicamente como contraste y no se
	utiliza para rellenar vacíos. De forma análoga, $HUMCAP$ se construye
	íntegramente con la serie de escolaridad del PNUD; el índice \texttt{hc} de
	PWT se conserva solo como referencia de validación y no completa las 27
	observaciones faltantes.

	\subsection{Disponibilidad y muestra efectiva}
	\label{subsec:cobertura-panel-maestro}

	La grilla maestra conserva las 1.430 llaves país--año aun cuando alguna
	variable no esté disponible. A partir de esa grilla, las estimaciones de
	$ECI$ y $DIVX$ utilizan una muestra común y fija de 1.044 observaciones,
	correspondientes a 49 países y 23 años efectivos dentro del intervalo
	1996--2021. El panel es desbalanceado: además de las ausencias específicas por
	país, la muestra común no contiene observaciones para 1997, 1999 y 2001. Por
	tanto, el intervalo calendario no equivale a una serie continua de 26 años.

	La muestra econométrica se obtiene por la intersección de las variables
	dependientes y los regresores requeridos, sin imputaciones ni interpolaciones
	y sin alterar la selección inicial de 55 economías definida mediante
	$DRES\geq20\%$ en 1990--1995. Esta regla permite que M1, M2 y M3 sean
	comparables y evita que las diferencias entre especificaciones respondan a
	cambios en la composición muestral.
	La cobertura individual, la composición de la muestra conjunta y los
	estadísticos descriptivos se presentan exclusivamente en el
	\autoref{chap:datos-variables}, para no duplicar resultados empíricos dentro
	del capítulo metodológico.

	\section{Especificación del modelo econométrico}
	\label{sec:modelo-econometrico}


	La especificación econométrica adoptada se presenta en la \autoref{eq:modelo_base} y se ilustra conceptualmente en la \autoref{fig:modelo-econometrico}. Esta propuesta desplaza el foco analítico desde la dependencia de recursos naturales como variable de resultado hacia la transformación estructural externa, aproximada principalmente mediante la complejidad económica.

	En este marco, la transformación estructural se aproxima desde su dimensión externa. La variable dependiente principal es la complejidad económica ($ECI$), que refleja la sofisticación de la canasta exportadora, y se considera como resultado complementario la diversificación exportadora ($DIVX=1-HHI$), que captura el grado de dispersión de dicha canasta.

	$DRES$ se utiliza exclusivamente para seleccionar la muestra y no interviene
	como variable dependiente ni como regresor. Por su parte, $RENTS$ aproxima la
	intensidad macroeconómica de las rentas extractivas y constituye uno de los
	componentes explicativos del modelo, no el resultado ni el único centro de la
	tesis. Su interacción con $INST$ permite evaluar si la asociación condicional
	entre las rentas extractivas y la transformación estructural varía con la
	calidad institucional; no identifica por sí misma un mecanismo causal de
	mediación.

	En términos analíticos, el modelo se organiza en seis canales complementarios:
	(i) rentas e instituciones, capturado por $RENTS$, $INST$ y su interacción;
	(ii) rentas extractivas específicas por habitante, medidas mediante $OILPC$,
	$GASPC$ y $COALPC$; (iii) estructura exportadora, representada por $HHI$,
	$PEXP$ y $FEXP$; (iv) estabilidad macroeconómica, aproximada mediante $VOL$ y
	$RER$; (v) capacidades productivas, reflejadas por $HUMCAP$, $INNOV$ y $NET$;
	y (vi) controles económicos y financieros, integrados por
	$\log(GDPPC_{it})$, $GOVCONS$ y $FIN$. $OILPC$, $GASPC$ y $COALPC$ son
	medidas monetarias de rentas reales per cápita y entran en las estimaciones
	mediante la transformación $\ln(1+x)$: no representan precios
	internacionales ni dotaciones físicas de recursos. Asimismo, $VOL$ es la
	desviación estándar móvil de cinco años de las variaciones logarítmicas del
	índice CTOT del FMI con ponderaciones fijas, y no un indicador de volatilidad
	financiera o del precio internacional del petróleo.

	La especificación econométrica de referencia se define como:

	\begin{equation}
		\label{eq:modelo_base}
		\begin{aligned}
			ECI_{it} = \alpha
			&+ \beta_1 RENTS_{it}
			+ \beta_2 INST_{it}
			+ \beta_3(RENTS_{it} \times INST_{it}) \\
			&+ \beta_4 \ln(1+OILPC_{it})
			+ \beta_5 \ln(1+GASPC_{it})
			+ \beta_6 \ln(1+COALPC_{it}) \\
			&+ \beta_7 HHI_{it}
			+ \beta_8 PEXP_{it}
			+ \beta_9 FEXP_{it} \\
			&+ \beta_{10} VOL_{it}
			+ \beta_{11} RER_{it} \\
			&+ \beta_{12} HUMCAP_{it}
			+ \beta_{13} INNOV_{it}
			+ \beta_{14} NET_{it} \\
			&+ \beta_{15}\log(GDPPC_{it})
			+ \beta_{16} GOVCONS_{it}
			+ \beta_{17} FIN_{it} \\
			&+ \mu_i + \lambda_t + \varepsilon_{it}
		\end{aligned}
	\end{equation}

	donde $ECI_{it}$ representa la complejidad económica del país $i$ en el período $t$, $\mu_i$ corresponde a efectos fijos por país, $\lambda_t$ a efectos temporales comunes y $\varepsilon_{it}$ al término de error. Esta especificación constituye el modelo principal de la tesis, dado que el ECI captura la sofisticación de la canasta exportadora y las capacidades productivas reveladas.

	La especificación incorpora de manera explícita las rentas extractivas
	específicas por habitante ($OILPC_{it}$, $GASPC_{it}$, $COALPC_{it}$), la
	estructura exportadora ($HHI_{it}$, $PEXP_{it}$, $FEXP_{it}$), las condiciones
	macroeconómicas ($VOL_{it}$, $RER_{it}$), las capacidades productivas
	($HUMCAP_{it}$, $INNOV_{it}$, $NET_{it}$) y los controles económicos y
	financieros ($\log(GDPPC_{it})$, $GOVCONS_{it}$, $FIN_{it}$). En esta
	especificación, $\log(GDPPC_{it})$ se interpreta explícitamente como proxy del
	nivel de desarrollo económico, y no como el PIB per cápita en niveles,
	mientras que $GOVCONS_{it}$ controla por el tamaño relativo del consumo
	público. Este conjunto de variables es consistente con el esquema conceptual
	presentado en la \autoref{fig:modelo-econometrico}.

	Desde el punto de vista teórico, se espera que una mayor concentración exportadora ($HHI$), así como condiciones macroeconómicas adversas ---particularmente la volatilidad de los términos de intercambio de commodities y la apreciación del tipo de cambio real---, afecten negativamente la complejidad económica. Por el contrario, mayores niveles de capital humano e innovación deberían favorecer la acumulación de capacidades productivas y la diversificación exportadora.

	En relación con el canal institucional, la asociación entre las rentas
	extractivas ($RENTS$) y la transformación estructural puede variar con la
	calidad institucional. La hipótesis plantea que una mayor intensidad de rentas
	puede asociarse con patrones de especialización en entornos institucionales
	débiles y con trayectorias productivas distintas en contextos institucionales
	más sólidos.

	La interacción $RENTS_{it} \times INST_{it}$ permite estimar asociaciones
	diferenciadas según el entorno institucional. Esta es la única interacción
	constitutiva incluida en las dos ecuaciones centrales del estudio.

	Como resultado complementario, la tesis estima una segunda especificación en la que la variable dependiente es la diversificación exportadora, definida como $DIVX_{it}=1-HHI_{it}$. Esta especificación permite evaluar si los mecanismos asociados a la complejidad económica también se relacionan con una menor concentración de la canasta exportadora:

	\begin{equation}
		\label{eq:modelo_divx}
		\begin{aligned}
			DIVX_{it} = \alpha
			&+ \delta_1 RENTS_{it}
			+ \delta_2 INST_{it}
			+ \delta_3(RENTS_{it} \times INST_{it}) \\
			&+ \delta_4 \ln(1+OILPC_{it})
			+ \delta_5 \ln(1+GASPC_{it})
			+ \delta_6 \ln(1+COALPC_{it}) \\
			&+ \delta_7 PEXP_{it}
			+ \delta_8 FEXP_{it} \\
			&+ \delta_9 VOL_{it}
			+ \delta_{10} RER_{it} \\
			&+ \delta_{11} HUMCAP_{it}
			+ \delta_{12} INNOV_{it}
			+ \delta_{13} NET_{it} \\
			&+ \delta_{14}\log(GDPPC_{it})
			+ \delta_{15} GOVCONS_{it}
			+ \delta_{16} FIN_{it} \\
			&+ \mu_i + \lambda_t + u_{it}
		\end{aligned}
	\end{equation}

	En la \autoref{eq:modelo_divx}, $HHI_{it}$ se excluye del conjunto de regresores porque $DIVX_{it}$ se construye directamente como su complemento. Esta es la única exclusión respecto del canal estructural: $PEXP_{it}$ y $FEXP_{it}$ permanecen como regresores y capturan, respectivamente, la especialización primaria no energética y la especialización en combustibles. De este modo, se evita introducir una relación mecánica entre la variable dependiente y una covariable explicativa sin alterar los demás canales del modelo. La interpretación de esta especificación es complementaria a la del modelo principal: mientras el ECI mide sofisticación y capacidades productivas reveladas, $DIVX$ captura la dispersión de la canasta exportadora.

	\subsection{Arquitectura de especificaciones M1, M2 y M3}
	\label{subsec:arquitectura-modelos}

	Para cada variable dependiente se estiman tres especificaciones sobre la misma
	muestra común. M1 resume el bloque de estructura extractiva y exportadora:
	incorpora $RENTS$, $INST$, su interacción, las tres rentas específicas per
	cápita, $PEXP$, $FEXP$ y $\log(GDPPC)$. En el modelo de $ECI$ también incluye
	$HHI$, mientras que en el modelo de $DIVX$ se excluye por construcción.

	M2 representa el bloque de capacidades y estabilidad. Incluye $RENTS$,
	$INST$, su interacción, $VOL$, $RER$, $HUMCAP$, $INNOV$, $NET$,
	$\log(GDPPC)$, $GOVCONS$ y $FIN$. M1 y M2 son especificaciones temáticas
	paralelas: ninguna constituye una ampliación secuencial de la otra y ambas se
	encuentran anidadas en M3.

	M3 reúne simultáneamente los seis canales descritos y corresponde a las
	ecuaciones completas de las \autoref{eq:modelo_base} y
	\autoref{eq:modelo_divx}. Por esta razón, constituye la especificación
	principal para la interpretación sustantiva y para el protocolo integral de
	diagnósticos y sensibilidad. M1 y M2 permiten examinar la estabilidad temática
	de los patrones, pero no sustituyen a M3 ni se seleccionan en función de la
	significancia obtenida.

	Todas las especificaciones incluyen efectos fijos por país y por año y
	errores estándar agrupados por país. Como extensión complementaria, la
	desagregación predefinida de $RENTS$ separa las rentas de petróleo y gas de las
	rentas de minería y carbón. Este ejercicio no modifica el centro analítico en
	$ECI$ y $DIVX$ ni se interpreta como una nueva variable de resultado.

	%%%%%%%%

	La \autoref{tab:signo} presenta las expectativas teóricas de signo para las variables incluidas en la especificación econométrica.

	\renewcommand{\arraystretch}{1.2}
	\setlength{\tabcolsep}{4pt}

	\begin{longtable}{|
			>{\raggedright\arraybackslash}p{0.20\textwidth}|
			>{\centering\arraybackslash}p{0.12\textwidth}|
			>{\raggedright\arraybackslash}p{0.58\textwidth}|}

		\caption[Expectativas teóricas de signo]
		{Expectativas teóricas de signo de las variables utilizadas en el modelo econométrico de transformación estructural.}

		\label{tab:signo} \\

		\hline

		Variable
		& Signo esperado
		& Justificación teórica \\

		\hline
		\endfirsthead

		\multicolumn{3}{c}
		{{\tablename\ \thetable{} -- continuación}} \\

		\hline

		Variable
		& Signo esperado
		& Justificación teórica \\

		\hline
		\endhead

		\hline
		\multicolumn{3}{r}{\textit{Continúa en la página siguiente}} \\
		\endfoot

		\hline

		\multicolumn{3}{p{\textwidth}}{
			\footnotesize
			\textit{Nota: Los signos esperados corresponden a las hipótesis teóricas planteadas para el modelo de transformación estructural externa. La variable dependiente principal es el Índice de Complejidad Económica (ECI), y la diversificación exportadora ($DIVX=1-HHI$) se utiliza como resultado complementario. Los signos pueden diferir entre ambas variables dependientes.}
		} \\

		\endlastfoot

		RENTS
		& (-)
		& Una mayor participación de las rentas de petróleo, gas natural, carbón y minerales en el PIB se asocia con menores niveles de transformación estructural debido a efectos de especialización extractiva, desplazamiento de sectores transables y dinámicas de enclave. \\
		\hline

		INST
		& (+)
		& Una mayor calidad institucional favorece la estabilidad de reglas, la asignación eficiente de recursos y la coordinación de políticas necesarias para la transformación productiva. \\
		\hline

		RENTS $\times$ INST
		& (+)
		& Instituciones de mayor calidad pueden atenuar la asociación adversa entre una mayor intensidad de rentas extractivas y la transformación estructural, y facilitar su conversión en capacidades productivas. \\
		\hline

		OILPC
		& (-)
		& Una mayor renta petrolera per cápita puede reforzar la especialización extractiva y desplazar actividades transables más complejas cuando no existen mecanismos eficaces de diversificación.
		\\
		\hline

		GASPC
		& (-)
		& Una mayor renta gasífera real por habitante puede profundizar patrones de enclave y dependencia de rentas, limitando la diversificación productiva.
		\\
		\hline

		COALPC
		& (-)
		& Una mayor renta carbonífera real por habitante tiende a asociarse con especialización en actividades extractivas de bajo encadenamiento local y menor transformación estructural.
		\\
		\hline

		HHI$^{*}$
		& (-)
		& Un mayor nivel de concentración productiva y exportadora indica especialización en pocos sectores o productos, lo que suele asociarse con menor diversificación y menor transformación estructural. Se incluye únicamente en las especificaciones donde $ECI$ es la variable dependiente. \\
		\hline

		PEXP
		& (-)
		& Una mayor participación de exportaciones primarias no energéticas en la canasta exportadora refleja especialización en bienes de bajo procesamiento y menores oportunidades de escalamiento productivo.
		\\
		\hline

		FEXP
		& (-)
		& Una mayor participación de combustibles en las exportaciones suele reforzar la dependencia de recursos energéticos y reducir los incentivos a diversificar la estructura productiva.
		\\
		\hline

		VOL
		& (-)
		& Una mayor volatilidad de los términos de intercambio de commodities incrementa la incertidumbre sobre los ingresos externos, puede desalentar la inversión de largo plazo y dificulta la acumulación sostenida de capacidades productivas. \\
		\hline

		RER (apreciación)
		& (-)
		& La apreciación del tipo de cambio real reduce la competitividad de sectores transables distintos del extractivo y limita la expansión de actividades manufactureras y complejas. \\
		\hline

		HUMCAP
		& (+)
		& Un mayor nivel de capital humano facilita la adopción tecnológica, el aprendizaje productivo y la diversificación hacia actividades de mayor complejidad. \\
		\hline

		INNOV
		& (+)
		& La innovación impulsa procesos de escalamiento productivo, difusión tecnológica y expansión hacia actividades de mayor sofisticación económica. \\
		\hline

		NET
		& (+)
		& Una mayor conectividad digital favorece la difusión de conocimiento, la coordinación productiva y la incorporación de tecnologías complementarias a la diversificación estructural.
		\\
		\hline

		$\log(GDPPC_{it})$
		& (+)
		& Un mayor valor de $\log(GDPPC_{it})$, utilizado como proxy del nivel de desarrollo económico, suele asociarse con estructuras productivas más diversificadas, mayores capacidades tecnológicas y mayor complejidad económica. \\
		\hline

		GOVCONS
		& ($\pm$)
		& El consumo final del gobierno como porcentaje del PIB se incorpora como control. Su asociación con la transformación estructural es teóricamente ambigua, pues puede reflejar provisión de capacidades públicas o desplazamiento de actividad privada. \\
		\hline

		FIN
		& (+)
		& Una mayor profundidad del crédito privado puede facilitar el financiamiento de actividades productivas e innovadoras, aunque el indicador no identifica por sí solo la asignación ni la calidad de ese crédito. \\
		\hline

	\end{longtable}
	\thesissource{Elaboración propia a partir del marco teórico y de las hipótesis de la investigación.}


	Las estimaciones se realizan sobre la muestra común desbalanceada de 1.044
	observaciones, 49 países y 23 años efectivos. M1, M2 y M3 incluyen efectos
	fijos por país y por año, y la inferencia principal utiliza errores estándar
	agrupados por país. Los modelos con efectos aleatorios no se emplean como
	referencia sustantiva.



	\begin{figure}[htbp]
		\centering
		\begin{tikzpicture}[
			scale=0.9,
			transform shape,
			font=\small,
			var/.style={anchor=west},
			channel/.style={anchor=west, font=\bfseries},
			dep/.style={
				draw,
				rectangle,
				rounded corners,
				minimum width=5.5cm,
				minimum height=1.5cm,
				align=center
			},
			arrow/.style={->, thick},
			node distance=0.2cm
			]

			% -------------------------------
			% CANAL INSTITUCIONAL
			% -------------------------------
			\node[channel] (c1) {Canal institucional};
			\node[var, below=of c1] (res) {RENTS};
			\node[var, below=of res] (inst) {INST};
			\node[var, below=of inst] (inter) {$RENTS \times INST$};

			% -------------------------------
			% RENTAS EXTRACTIVAS ESPECÍFICAS
			% -------------------------------
			\node[channel, below=0.5cm of inter] (cA) {Rentas per cápita};
			\node[var, below=of cA] (oil) {$\ln(1+OILPC)$};
			\node[var, below=of oil] (gas) {$\ln(1+GASPC)$};
			\node[var, below=of gas] (coal) {$\ln(1+COALPC)$};

			% -------------------------------
			% CANAL ESTRUCTURAL
			% -------------------------------
			\node[channel, below=0.5cm of coal] (c2) {Canal estructural};
			\node[var, below=of c2] (hhi) {HHI$^{*}$};
			\node[var, below=of hhi] (pexp) {PEXP};
			\node[var, below=of pexp] (fexp) {FEXP};

			% -------------------------------
			% CANAL MACRO
			% -------------------------------
			\node[channel, below=0.5cm of fexp] (c3) {Canal macroeconómico};
			\node[var, below=of c3] (vol) {VOL};
			\node[var, below=of vol] (rer) {RER};

			% -------------------------------
			% CAPACIDADES
			% -------------------------------
			\node[channel, below=0.5cm of rer] (c4) {Capacidades productivas};
			\node[var, below=of c4] (hum) {HUMCAP};
			\node[var, below=of hum] (inn) {INNOV};
			\node[var, below=of inn] (net) {NET};

			% -------------------------------
			% CONTROL
			% -------------------------------
			\node[channel, below=0.5cm of net] (c5) {Control};
			\node[var, below=of c5] (gdp) {$\log(GDPPC_{it})$};

			% -------------------------------
			% FINANCIERO
			% -------------------------------
			\node[channel, below=0.5cm of gdp] (c6) {Controles fiscales y financieros};
			\node[var, below=of c6] (govcons) {GOVCONS};
			\node[var, below=of govcons] (fin) {FIN};

			% -------------------------------
			% VARIABLE DEPENDIENTE
			% -------------------------------
			\node[dep, right=5cm of vol] (depvar) {
				\textbf{Transformación estructural}\\
				\textit{ECI}\\
				\textit{DIVX = 1 - HHI}
			};

			% -------------------------------
			% FLECHAS
			% -------------------------------
			\draw[arrow] (res.east)   -- node[above] {$-$} (depvar.west);
			\draw[arrow] (inst.east)  -- node[above] {$+$} (depvar.west);
			\draw[arrow] (inter.east) -- node[above] {$+$} (depvar.west);

			\draw[arrow] (oil.east)  -- node[above] {$-$} (depvar.west);
			\draw[arrow] (gas.east)  -- node[above] {$-$} (depvar.west);
			\draw[arrow] (coal.east) -- node[above] {$-$} (depvar.west);

			\draw[arrow] (hhi.east)  -- node[above] {$-$} (depvar.west);
			\draw[arrow] (pexp.east) -- node[above] {$-$} (depvar.west);
			\draw[arrow] (fexp.east) -- node[above] {$-$} (depvar.west);

			\draw[arrow] (vol.east)  -- node[above] {$-$} (depvar.west);
			\draw[arrow] (rer.east)  -- node[above] {$-$} (depvar.west);

			\draw[arrow] (hum.east)  -- node[above] {$+$} (depvar.west);
			\draw[arrow] (inn.east)  -- node[above] {$+$} (depvar.west);
			\draw[arrow] (net.east)  -- node[above] {$+$} (depvar.west);

			\draw[arrow] (gdp.east)  -- node[above] {$+$} (depvar.west);

			\draw[arrow] (govcons.east) -- node[above] {$\pm$} (depvar.west);
			\draw[arrow] (fin.east)  -- node[above] {$+$} (depvar.west);

		\end{tikzpicture}

		\caption[Modelo econométrico de transformación estructural externa]{Modelo econométrico de transformación estructural externa. Esquema conceptual organizado por canales teóricos, incorporando rentas extractivas reales per cápita, estructura exportadora, estabilidad macroeconómica, capacidades productivas y diversificación exportadora como resultado complementario.\protect\\[0.3em]\textit{Nota:} $HHI$ se incluye en el modelo donde $ECI$ es la variable dependiente; en el modelo complementario con $DIVX=1-HHI$, se excluye del conjunto de regresores.}
		\label{fig:modelo-econometrico}
		\caption*{\footnotesize \textit{Fuente: Elaboración propia.}}
	\end{figure}


	\section{Procedimientos de análisis econométrico}
	\label{sec:analisis-econometrico}

	El análisis empírico se organiza en seis etapas. En primer lugar, se describe
	la selección inicial de 55 economías mediante $DRES\geq20\%$ en 1990--1995 y
	se examinan la cobertura temporal, los vacíos y las distribuciones de las
	variables, sin imputar ni interpolar observaciones.

	En segundo lugar, se fija la muestra econométrica común de 1.044 observaciones,
	49 países y 23 años efectivos del período 1996--2021. Esta misma muestra se
	mantiene en M1, M2 y M3 para $ECI$ y $DIVX$, con la única exclusión sustantiva
	de $HHI$ en las ecuaciones de $DIVX$.

	En tercer lugar, se estiman M1 y M2 como especificaciones temáticas paralelas y
	M3 como especificación completa principal. En todos los casos se incorporan
	efectos fijos por país y por año y se calculan errores estándar agrupados por
	país. La interacción $RENTS_{it}\times INST_{it}$ se incluye junto con sus
	términos constitutivos.

	En cuarto lugar, se realiza un diagnóstico descriptivo de persistencia mediante
	coeficientes autorregresivos AR(1), sin convertirlos en categorías mecánicas,
	y se aplican pruebas Fisher--ADF como información complementaria sobre el
	comportamiento temporal de las series. Sus resultados no se utilizan para
	clasificar automáticamente las variables como $I(0)$ o $I(1)$, transformar
	las series ni seleccionar el estimador. Sobre los residuos de M3 se aplican la
	prueba de Wald modificada para heterocedasticidad, la prueba de Wooldridge para
	autocorrelación y la prueba de Pesaran CD para dependencia transversal; la
	colinealidad se examina mediante factores de inflación de la varianza en la
	transformación \textit{within}. Estas pruebas se interpretan como diagnósticos y no
	como reglas automáticas de reemplazo de la especificación principal.

	En quinto lugar, la inferencia y la estabilidad de M3 se evalúan mediante
	\textit{Wild Cluster Bootstrap} con 9.999 replicaciones
	(\textcite{CameronGelbachMiller2008};
	\textcite{Roodman2019WildBootstrap}), exclusión iterativa de países
	(\textit{leave-one-out}), análisis de observaciones influyentes y cálculo de
	efectos marginales de $RENTS$ dentro del soporte observado de $INST$. También
	se estiman especificaciones con tendencias lineales por país. El año 2014 se
	incorpora únicamente como sensibilidad temporal preespecificada mediante
	interacciones en la muestra completa y pruebas conjuntas; no se divide el
	panel en regresiones separadas antes y después de ese año ni se interpreta
	como un quiebre causal.

	En sexto lugar, se estima la desagregación complementaria que distingue las
	rentas de petróleo y gas de las rentas de minería y carbón. Este ejercicio
	compara asociaciones sectoriales sin modificar la selección muestral ni
	atribuir carácter causal a las diferencias observadas.

	En conjunto, el diseño identifica asociaciones condicionales dentro de los
	países, netas de heterogeneidad invariable y de shocks anuales comunes. No
	permite afirmar que las variables explicativas causen enfermedad holandesa,
	destrucción de capacidades o diversificación. Los diagnósticos y ejercicios de
	sensibilidad delimitan la estabilidad de los resultados, pero no eliminan la
	endogeneidad propia de un diseño observacional, comparativo y no experimental.


	\section{Métodos econométricos descartados y fundamentación metodológica}
	\label{sec:metodos-descartados}

	La elección de la estrategia empírica responde a la pregunta de investigación,
	a la estructura observacional del panel ($N=49$ países y 23 años efectivos,
	con un promedio de 21,3 observaciones por país) y a la ausencia de una fuente
	exógena de variación que permita sostener identificación causal. Las
	alternativas se descartan por razones de diseño y no a partir de comparaciones
	retrospectivas de significancia estadística.

	\subsection{Delimitación metodológica, diagnósticos y justificación de la especificación}
	\label{subsec:delimitacion-metodologica}

	La estrategia se circunscribe a asociaciones condicionadas estimadas mediante
	efectos fijos por país y año. Esta delimitación no supone que otros métodos
	carezcan de utilidad en paneles macroeconómicos; indica que, con la pregunta,
	las series y la estructura muestral de esta investigación, no aportan una
	identificación adicional suficientemente defendible. La inclusión o exclusión
	de cada extensión se definió antes de interpretar los coeficientes principales.

	\noindent\textit{Diagnósticos e inferencia:} Los residuos de M3 presentan
	heterocedasticidad entre países y autocorrelación de primer orden. Por ello,
	la inferencia principal utiliza errores estándar agrupados por país, que
	permiten dependencia arbitraria dentro de cada conglomerado. Como contraste
	para el número moderado de 49 conglomerados, los términos focales se evalúan
	además mediante \textit{wild cluster bootstrap} con 9.999 replicaciones
	\parencite{Cameron2008, Roodman2019}.
	La prueba de Pesaran CD aplicada a residuos con efectos fijos anuales no
	rechaza la independencia transversal ni para $ECI$ ($p=0{,}663$) ni para
	$DIVX$ ($p=0{,}442$) \parencite{Pesaran2004}. En consecuencia,
	Driscoll--Kraay no sustituye la inferencia agrupada ni se activa
	automáticamente como sensibilidad; su uso requeriría una justificación,
	especificación y regla de rezagos separadas
	\parencite{Driscoll1998}.

	\noindent\textit{Cointegración y modelos de corrección de errores:} No se estiman
	pruebas de cointegración en panel ni modelos de corrección de errores. Las
	pruebas Fisher--ADF entregan evidencia mixta entre variables y variantes, los
	huecos internos del panel limitan la aplicación uniforme de otras pruebas y
	no se preespecificó un vector de largo plazo que reproduzca exactamente la
	ecuación M3 completa. Bajo estas condiciones, una prueba aplicada a un
	subconjunto de regresores no justificaría interpretar la relación estimada
	como un equilibrio de largo plazo \parencite{Pedroni1999, Kao1999}. Los efectos fijos no eliminan por sí solos los problemas
	de integración; por esa razón, la tesis se limita explícitamente a
	asociaciones contemporáneas condicionadas y no formula afirmaciones de
	equilibrio de largo plazo.

	\noindent\textit{Filtro HP y proyecciones locales:} Tampoco se utiliza el filtro
	Hodrick--Prescott. Una descomposición mecánica entre tendencia y ciclo no
	resuelve la endogeneidad de un diseño observacional ni identifica una relación
	estructural de largo plazo; además, puede introducir dinámicas artificiales y
	distorsiones en los extremos de la muestra \parencite{Hamilton2018}. Filtrar
	solo $ECI$ o $DIVX$ mientras los regresores permanecen en niveles agravaría la
	asimetría de la especificación. Por su parte, las proyecciones locales pueden
	describir asociaciones dinámicas, pero no permiten denominar sus coeficientes
	respuestas a impulsos causales cuando no existe un \textit{shock} identificado
	ni una estrategia de identificación externa \parencite{Jorda2005}. Por ello,
	no se incorporan como evidencia del TFM.

	\subsection{Métodos adicionales descartados}

	\begin{enumerate}
		\item \textit{Panel dinámico mediante GMM}: los estimadores GMM pueden
		ser vulnerables a la proliferación de instrumentos y a pruebas de validez
		poco informativas cuando la dimensión temporal no es corta
		\parencite{Roodman2009}. Además, la pregunta central no
		requiere incluir una variable dependiente rezagada como componente del
		modelo principal.
		\item \textit{Variables instrumentales y 2SLS}: no se dispone de un
		instrumento externo que satisfaga de manera defendible la relevancia y la
		restricción de exclusión para $ECI$ y $DIVX$
		\parencite{Angrist1996, Andrews2019}. Los
		rezagos internos no resuelven por sí mismos la endogeneidad asociada a
		shocks persistentes y decisiones de política.
		\item \textit{Sistemas dinámicos multivariados}: un sistema que trate como
		endógeno el conjunto completo de variables no corresponde al objetivo de
		contrastar los canales teóricos definidos y exigiría una parametrización
		desproporcionada frente al tamaño efectivo del panel. Por ello, no se adopta
		como evidencia principal.
		\item \textit{Diseños de tratamiento, censura o selección}: el umbral de
		$DRES$ define ex ante la población analítica, pero no constituye una regla
		de asignación cuasiexperimental. Tampoco existe censura de las variables
		dependientes que justifique modelos de respuesta limitada. En consecuencia,
		estos diseños no solucionan la pregunta ni las limitaciones de identificación
		del estudio.
	\end{enumerate}


% =====================================
% FIN DE TU CONTENIDO
% =====================================

	\thesisendchap{metodologia}

\end{document}
